\section*{Leitidee}

Die reellen Zahlen reichen nicht aus, um jede algebraische Gleichung zu
l�sen. Bereits die einfache Gleichung
\[
x^2=-1
\]
besitzt keine reelle L�sung. Die Einf�hrung der \textbf{komplexen Zahlen}
(\emph{complex numbers}) erweitert den bisherigen Zahlenbereich so, dass auch
solche Gleichungen sinnvoll behandelt werden k�nnen.

Komplexe Zahlen verbinden dabei mehrere mathematische Perspektiven. In
\textbf{kartesischer Darstellung} erscheinen sie als Zahlen der Form
\[
	z=x+ i \, y
\]
und k�nnen geometrisch als Punkte oder Zeiger in der komplexen Ebene
interpretiert werden. Dadurch erhalten auch algebraische Operationen eine
geometrische Bedeutung.

Eine zweite Darstellung entsteht durch Betrag und Winkel. Diese
\textbf{Polarform} beschreibt eine komplexe Zahl durch ihre L�nge und ihre
Richtung. Mithilfe der komplexen Exponentialfunktion f�hrt sie auf die
Darstellung
\[
	z=r\,e^{i\,\varphi}.
\]
Insbesondere Multiplikation, Division, Potenzieren und Wurzelziehen erhalten
dadurch eine besonders anschauliche und zugleich rechnerisch g�nstige Form.

Die Verbindung zur Exponentialfunktion f�hrt schlie�lich unmittelbar zu
periodischen Vorg�ngen. Die komplexe Exponentialfunktion beschreibt eine
Drehbewegung in der komplexen Ebene und bildet damit die Grundlage der
\textbf{Zeigerdarstellung} harmonischer Schwingungen. Amplitude und Phase
werden dabei durch Betrag und Winkel einer komplexen Zahl beschrieben.
�berlagerung, Ableitung und Integration harmonischer Schwingungen k�nnen auf
diese Weise auf algebraische Operationen mit komplexen Zahlen zur�ckgef�hrt
werden.

Die Verbindung von algebraischer Erweiterung, geometrischer Interpretation
und analytischer Darstellung macht die komplexen Zahlen zu einem
grundlegenden Werkzeug f�r die Beschreibung periodischer und dynamischer
Prozesse.

